(a) O grupo adimensional proposto é
\[
\Pi = \frac{P}{E L^2},
\]
que é sem dimensão (força dividida por módulo vezes área equivalente de escala).

(b) Pela igualdade de grupos entre modelo (m) e protótipo (p):
\[
\frac{P_m}{E_m L_m^2}=\frac{P_p}{E_p L_p^2}
\Rightarrow
P_m=P_p\frac{E_m}{E_p}\left(\frac{L_m}{L_p}\right)^2.
\]
Com $P_p=9000$ N, $L_m=0{,}20$ m e $L_p=3{,}6$ m:
\[
P_m=9000\,\frac{E_m}{E_p}\left(\frac{0{,}20}{3{,}6}\right)^2.
\]
Substituindo os módulos de aço e madeira adotados, obtém-se o valor numérico final.

